\section{Verification of the Core Model}
\label{sec:verification}

Across five theories (\Cref{tab:theories}) we discharge 87 lemma checks over 61
distinct properties. The concrete theory \code{multihop.spthy} contributes 43
lemmas: 33 all-traces and 10 exists-trace witnesses, of which 5 all-traces
lemmas are \code{[reuse]} helpers rather than protocol goals. The $N$-hop
theory \code{Modif.spthy} re-proves 26 of these over an abstract routing layer
and adds 5 new ones, for 31 in total. The remaining 13 come from the three
timing theories.

\Cref{fig:multihop} shows the flow the concrete theory models, with each stage
annotated by the lemmas that constrain it.

\begin{figure*}[t]
\centering
\begin{tikzpicture}[
  font=\sffamily,
  x=1cm,y=1cm,
  actor/.style ={rounded corners=2pt,fill=dgreen,text=white,
                 font=\sffamily\scriptsize\bfseries,inner xsep=6pt,inner ysep=3pt},
  life/.style  ={draw=cRule,densely dashed,line width=.4pt},
  msg/.style   ={draw=mgreen,line width=.8pt,-{Stealth[length=4.5pt,width=3.5pt]}},
  msgd/.style  ={draw=mgreen,line width=.8pt,densely dashed,
                 -{Stealth[length=4.5pt,width=3.5pt]}},
  omsg/.style  ={draw=cOrange,line width=.8pt,densely dashed,
                 -{Stealth[length=4.5pt,width=3.5pt]}},
  rmsg/.style  ={draw=cRed,line width=.8pt,
                 dash pattern=on 3pt off 1.5pt on .6pt off 1.5pt,
                 -{Stealth[length=4.5pt,width=3.5pt]}},
  lbl/.style   ={font=\sffamily\scriptsize,text=black,inner sep=2pt},
  bdg/.style   ={circle,minimum size=3.6mm,inner sep=0pt,
                 text=white,font=\sffamily\tiny\bfseries},
  band/.style  ={fill=lgreen,rounded corners=2pt},
  rule/.style  ={draw=cRule,line width=.4pt},
]

%================================================= geometry
\def\xA{1.35} \def\xB{4.15} \def\xC{6.95} \def\xD{9.75}
\def\xL{0.0}  \def\xR{10.70} \def\xSep{11.05} \def\xRail{11.35}
\def\xBdg{0.85}

%================================================= legend
\foreach \lx/\lc/\lt in {%
   0.05/mgreen/{honest success / safety proved},
   4.25/cOrange/{timeout / refund},
   6.55/cRed/{attack (reachable)}}{%
  \fill[\lc] (\lx,-0.12) rectangle ++(0.24,0.24);
  \node[anchor=west,font=\sffamily\scriptsize] at (\lx+0.32,0) {\lt};}
\node[anchor=west,font=\sffamily\tiny,text=cGray]
  at (0.05,-0.42) {circled numbers $=$ lemmas on the right, not message order};

%================================================= lifelines
\foreach \x/\n in {\xA/Alice,\xB/Bob,\xC/Carol,\xD/Dave}{%
  \node[actor] at (\x,-0.75) {\n};
  \draw[life] (\x,-1.05) -- (\x,-9.35);}

%================================================= round bands
\node[band,minimum width=.42cm,minimum height=2.70cm,
      anchor=north west] (r1) at (\xL,-1.40) {};
\node[rotate=90,font=\sffamily\tiny\bfseries,text=dgreen] at (r1.center)
      {ROUND 1 $\cdot$ LOCK};
\node[band,minimum width=.42cm,minimum height=1.90cm,
      anchor=north west] (r2) at (\xL,-4.45) {};
\node[rotate=90,font=\sffamily\tiny\bfseries,text=dgreen] at (r2.center)
      {ROUND 2 $\cdot$ RELEASE};

%================================================= round 1 : lock
\draw[msgd] (\xD,-1.65) -- node[lbl,above]
      {Dave samples $x$, sends invoice $y = H(x)$ for 3 to Alice} (\xA,-1.65);
\node[bdg,fill=mgreen] at (\xBdg,-1.65) {1};

\draw[msg] (\xA,-2.45) -- node[lbl,above] {HTLC($A\!\to\!B$, 5, $y$, $3t$)} (\xB,-2.45);
\node[bdg,fill=mgreen] at (\xBdg,-2.45) {2};
\draw[msg] (\xB,-3.10) -- node[lbl,above] {HTLC($B\!\to\!C$, 4, $y$, $2t$)} (\xC,-3.10);
\draw[msg] (\xC,-3.75) -- node[lbl,above] {HTLC($C\!\to\!D$, 3, $y$, $t$)}  (\xD,-3.75);

\draw[rule] (\xBdg,-4.25) -- (\xR,-4.25);

%================================================= round 2 : release
\draw[msg] (\xD,-4.80) -- node[lbl,above] {redeem($x$) --- Dave reveals $x$} (\xC,-4.80);
\node[bdg,fill=mgreen] at (\xBdg,-4.80) {3};
\draw[msg] (\xC,-5.45) -- node[lbl,above] {redeem($x$)} (\xB,-5.45);
\draw[msg] (\xB,-6.10) -- node[lbl,above] {redeem($x$)} (\xA,-6.10);

\node[bdg,fill=mgreen] at (\xBdg,-6.65) {4};
\node[anchor=west,font=\sffamily\scriptsize\bfseries,text=dgreen]
      at (\xA+0.45,-6.65)
      {$\to$ payment successful (each party paid; fee to each intermediary)};
\node[bdg,fill=mgreen] at (\xD+0.50,-6.65) {5};

%================================================= what can go wrong
\node[anchor=west,font=\sffamily\scriptsize\bfseries] at (\xBdg+0.35,-7.05)
      {WHAT CAN GO WRONG};
\draw[rule] (\xBdg+3.45,-7.05) -- (\xR,-7.05);

\draw[omsg] (\xD,-8.35) -- (\xA-0.1,-8.35);
\node[lbl,anchor=south west,text=cOrange,align=left] at (\xA,-8.27)
      {if $x$ never arrives, each HTLC refunds after its own timeout;\\
       staggering $t < 2t < 3t$ lets every intermediary redeem its\\
       inbound HTLC before its outbound one can expire.};
\node[bdg,fill=cOrange] at (\xBdg,-8.35) {6};

\draw[rmsg] (\xC+0.8,-9.20) -- (\xA-0.1,-9.20);
\node[lbl,anchor=south west,text=cRed,align=left] at (\xA,-9.12)
      {wormhole: Alice and Carol collude and relay $x$ off-path,\\
       skipping Bob --- the sender pays $f_1{+}f_2$ that no one earns.};
\node[bdg,fill=cRed] at (\xBdg,-9.20) {7};

%================================================= staggered timelocks
\node[bdg,fill=dgreen] at (\xL+0.22,-9.80) {8};
\node[anchor=west,font=\sffamily\scriptsize\bfseries] at (\xL+0.50,-9.80)
      {STAGGERED TIMELOCKS (block height $\to$)};

\def\yT{-10.45}
\fill[lgreen] (\xA,\yT+0.26) rectangle (\xA+2.2,\yT+0.42);
\fill[lgreen] (\xA,\yT+0.08) rectangle (\xA+4.4,\yT+0.24);
\fill[lgreen] (\xA,\yT-0.10) rectangle (\xA+6.6,\yT+0.06);
\draw[msg] (\xA,\yT-0.26) -- (\xA+7.55,\yT-0.26);
\node[lbl,anchor=west] at (\xA+7.65,\yT-0.26) {nested windows};

\foreach \dx/\tk in {2.2/{$t$ (C$\to$D)},4.4/{$2t$ (B$\to$C)},6.6/{$3t$ (A$\to$B)}}{%
  \draw[cRed,line width=.7pt] (\xA+\dx,\yT+0.46) -- (\xA+\dx,\yT-0.42);
  \node[font=\sffamily\scriptsize,anchor=north] at (\xA+\dx,\yT-0.50) {\tk};}

%================================================= lemma rail
\draw[rule] (\xSep,-0.10) -- (\xSep,\yT-0.60);
\node[anchor=north west,font=\sffamily\scriptsize\bfseries] at (\xRail,0.14)
      {LEMMAS $\cdot$ multihop.spthy};

\newcommand{\lemmarow}[5]{%
  \node[bdg,fill=#2] at (\xRail+0.13,#1) {#3};
  \node[anchor=north west,font=\sffamily\scriptsize\bfseries,text=#2,inner sep=0pt]
       at (\xRail+0.40,#1+0.11) {#4};
  \node[anchor=north west,font=\sffamily\tiny,inner sep=0pt,text width=5.95cm,
        align=left,inner ysep=1pt] at (\xRail+0.40,#1-0.09) {#5};}

\lemmarow{-1.00}{mgreen}{1}{Invoice \& preimage authentication}
  {Invoice\_Authenticates\_Settlement \md Preimage\_Secret\_Until\_Released\\
   Forged\_Invoice\_Requires\_Key\_Compromise \md Invoice\_Released\_Once \ck}

\lemmarow{-1.95}{mgreen}{2}{Lock causality (forward needs previous)}
  {Forward1\_Requires\_Offer \md Forward2\_Requires\_Forward1\\
   HTLC\_On\_Opened\_Channel \ck}

\lemmarow{-2.90}{mgreen}{3}{Release causality (redeem needs preimage)}
  {Fulfill\_Requires\_Forward2 \md Claim\_Requires\_Release\\
   Settle\_Requires\_Receiver\_Release \ck}

\lemmarow{-3.85}{mgreen}{4}{Atomicity + non-vacuity}
  {Payment\_Atomicity\_Under\_Liveness \ck\\
   Settle\_Excludes\_Sender\_Refund \ck\ (no refund after settle)\\
   Multihop\_Payment\_Possible \ck\ (witness)}

\lemmarow{-5.00}{mgreen}{5}{Value \& fee conservation}
  {EndToEnd\_Value\_Conservation \md Fee\_Conservation\_Hop1/2\\
   Fee\_Strictly\_Positive\_Hop1/2 \md Receiver\_Paid\_Invoice\_Amount \ck\\
   Fees\_Charged\_On\_Path\_Possible \ck\ (witness: both hops earn)}

\lemmarow{-6.15}{cOrange}{6}{Refund / griefing}
  {Refund\_Requires\_Timeout \md Loss\_Requires\_Inaction \ck\\
   Honest\_Intermediary\_Refunds \ck\ (reachable refund/DoS)}

\lemmarow{-7.10}{cRed}{7}{Wormhole fee theft}
  {Wormhole\_Fee\_Theft\_Reachable \ck\\
   Wormhole\_Steals\_Exactly\_The\_Fees \ck\ (theft = path fees)}

\lemmarow{-8.05}{dgreen}{8}{Intermediary safety (under liveness T2b)}
  {Intermediary\_Never\_Loses\_Under\_Liveness \ck\\
   Timeout\_Race\_Blocked \ck\ (race blocked with T2b)}

\node[anchor=north west,font=\sffamily\tiny,text width=6.2cm,align=left]
     at (\xRail-0.05,-9.05)
     {\ck\ = machine-checked in multihop.spthy (43/43 lemmas;
      the rail shows a selection).\\[2pt]
      green = all-traces safety \md orange = benign/reachable \md
      red = attack shown reachable (by design).};

\end{tikzpicture}
\caption{Multi-hop HTLC from Alice to Dave, annotated with lemmas from
  \texttt{multihop.spthy}. Circled numbers match the rail on the right
  (not message order). Green: proved safety; orange: reachable refund
  (griefing); red: reachable wormhole. Hop amounts $5\!\to\!4\!\to\!3$
  are the fees.}
\label{fig:multihop}
\end{figure*}


\subsection{Linear Facts and the Single-Spend Discipline}
\label{sec:discipline}

The central modelling lesson of this work concerns the choice between linear
and persistent facts. Recall from \Cref{sec:notation} that a persistent fact
\code{!F(t)} is checked by a rule but not consumed, whereas a linear fact
\code{F(t)} is removed from the state when the rule fires. This distinction is
standard; what is not obvious is which side of it an on-chain output belongs
on.

\begin{defn}[Single-spend discipline]
\label{def:singlespend}
A protocol resource that corresponds to \emph{one} on-chain output, spendable
at most once, must be modelled by a linear fact. A resource that corresponds to
\emph{knowledge}, which cannot be un-learned, may be modelled by a persistent
fact. Where a single event creates both --- a secret becomes known \emph{and}
an output becomes spendable --- the two must be separated into distinct facts.
\end{defn}

Lightning's revocation mechanism is exactly such an event, and it is a trap.
When a channel state is revoked, the counterparty learns the revocation secret
(knowledge, persistent) and simultaneously gains the right to spend one revoked
commitment output (a resource, linear). The natural encoding conflates them,
because the protocol description presents revocation as the disclosure of a
secret, and disclosure is what persistent facts are for.

\medbfparagraph{Instance 1: unbounded punishment}
Our first formulation gated the punishment rule on the persistent revocation
secret alone:

\begin{lstlisting}
rule Publish_Revoked_A:   // WRONG
    [ !ChannelRecord(ptr, $A, $B), !ValidOnChain(ptr),
      !RevokedSecret($A, n_old, sAold, oldCommitA) ]
    --[ PublishRevoked($A, n_old, oldCommitA), Cheat($A, n_old) ]->
    [ OnChainCheat($A, $B, n_old, oldCommitA),
      RevealedSecret($A, n_old, sAold, oldCommitA) ]
\end{lstlisting}

Every premise here is persistent, so the rule may fire arbitrarily often for a
single revoked state. The model admitted traces in which the honest party
confiscates the channel balance repeatedly on the strength of one act of
cheating. The defect is silent: every individual safety lemma about punishment
continues to hold. \code{No\_Punishment\_Without\_Cheating} is unaffected, since
punishment still requires a prior cheat. Only the \emph{multiplicity} is wrong,
and no single safety lemma in the original formulation constrained
multiplicity.

The repair separates the two roles. \code{Revoke\_Old\_Secret} now emits a
linear token alongside the persistent secret:

\begin{lstlisting}
    [ StateA(...), StateB(...),
      !RevokedSecret($A, n_old, sAold, oldCommitA),
      !RevokedSecret($B, n_old, sBold, oldCommitB),
      // linear one-shot spend of the revoked commitment output
      RevokedCommitmentUnspent(ptr, $A, $B, n_old, oldCommitA),
      RevokedCommitmentUnspent(ptr, $B, $A, n_old, oldCommitB) ]
\end{lstlisting}

and \code{Publish\_Revoked\_A} consumes that token, so the same commitment
cannot enter the punishment path twice:

\begin{lstlisting}
rule Publish_Revoked_A:   // FIXED
    [ !ChannelRecord(ptr, $A, $B), !ValidOnChain(ptr),
      !RevokedSecret($A, n_old, sAold, oldCommitA),
      RevokedCommitmentUnspent(ptr, $A, $B, n_old, oldCommitA) ]
    --[ PublishRevoked($A, n_old, oldCommitA), Cheat($A, n_old) ]->
    [ OnChainCheat($A, $B, n_old, oldCommitA),
      RevealedSecret($A, n_old, sAold, oldCommitA) ]
\end{lstlisting}

Knowledge of the secret is retained; only the spend is one-shot.

\medbfparagraph{Instance 2: HTLC slot reuse}
The same pathology appeared independently in the routing layer. A channel's
HTLC output was re-derivable from the persistent \code{!ChannelConnect} fact, so
one channel could carry unboundedly many concurrent HTLCs on the same output.
The repair is identical in shape: \code{Lock\_Funds\_And\_Open} mints a linear
\code{Free(ptr)} token, and every rule that adds an HTLC consumes it.

That two independent gaps in different layers had the same cause, and the same
one-line repair, is what makes \Cref{def:singlespend} worth stating as a
discipline rather than as two bug fixes.

\medbfparagraph{Re-verification}
Adding a linear premise shrinks a rule's enabling condition, so
$\mathrm{traces}(R') \subseteq \mathrm{traces}(R)$: every all-traces property is
preserved automatically, but exists-trace witnesses may be destroyed. We
therefore re-ran the full theory rather than appealing to monotonicity. All 43
lemmas re-verify. Notably the hardened $N$-hop theory verifies \emph{faster}
than the unhardened one ($134$\,s versus $265$\,s), because each additional
linear premise prunes the search space --- the opposite of the intuition that
had previously discouraged us from combining the two.

\subsection{Protocol Properties}

We state the main guarantees as theorems, each naming the \mtamarin\ lemma that
discharges it. Throughout, ``honest'' abbreviates the absence of any long-term
key compromise, which our model exposes as an explicit adversarial capability
(\code{Compromise\_Ltk}).

\begin{thrm}[Payment atomicity]
\label{thm:atomicity}
If both intermediaries forwarded, the receiver's HTLC was redeemed, and no key
was compromised, then the sender's HTLC on the same hash cannot be refunded.
\hfill(\code{Payment\_Atomicity\_Under\_Liveness})
\end{thrm}

\begin{thrm}[Intermediary safety]
\label{thm:intermediary}
If a forwarder's outgoing HTLC is redeemed and its incoming HTLC is refunded,
then its long-term key was compromised.
\hfill(\code{Intermediary\_Never\_Loses\_Under\_Liveness})
\end{thrm}

\Cref{thm:intermediary} is the contrapositive of the property that matters
operationally: an honest forwarder never pays out downstream without recovering
upstream. It holds only under the claim obligation T3, as
\Cref{sec:minimality} shows.

\begin{thrm}[Routing causality]
\label{thm:causality}
Every forward is preceded by a genuine upstream offer or forward with agreeing
amounts, unless the upstream party's key was compromised.
\hfill(\code{Forward1\_Requires\_Offer}, \code{Forward2\_Requires\_Forward1})
\end{thrm}

The compromise-free versions are proved separately as
\code{Forward1\_Requires\_Offer\_Honest} and
\code{Forward2\_Requires\_Forward1\_Honest}, which drop the disjunct rather than
carrying it.

\begin{thrm}[Value conservation]
\label{thm:value}
Over a complete linked three-hop payment, the sender locks strictly more than
the receiver's invoice amount.
\hfill(\code{EndToEnd\_Value\_Conservation})
\end{thrm}

The surplus is identified with the per-hop fees by
\code{Fee\_Conservation\_Hop1} and \code{Fee\_Conservation\_Hop2}, each of which
states $v_{\mathrm{in}} = v_{\mathrm{out}} + f$ at its hop. Taken together
these give the quantitative statement the wormhole analysis needs: value is
neither created nor destroyed along the path, and the difference between what
the sender pays and what the receiver receives is exactly $\sum_i f_i$.

\medbfparagraph{Non-vacuity}
Safety properties are vacuously true of a model that cannot execute. We
therefore accompany them with exists-trace witnesses.
\code{Multihop\_Payment\_Possible} exhibits a complete eight-event payment from
invoice to sender settlement. \code{Distinct\_Parties\_Configuration} exhibits
three channels between four pairwise-distinct parties with no compromise and no
dispute. \code{Fees\_Charged\_On\_Path\_Possible} exhibits two \emph{different}
intermediaries each charging a fee on the same payment.

\subsection{Timing Assumptions and Their Minimality}
\label{sec:minimality}

\Cref{sec:threat} introduced four timing restrictions. Two questions follow:
are they sound, and is any of them redundant?

\medbfparagraph{T2b is load-bearing}
With T2b active, the early-timeout race is unreachable
(\code{Timeout\_Race\_Blocked}). A minimal model with no such assumption
exhibits the race as a concrete trace (\code{timeout.spthy},
\code{Early\_Timeout\_Race}). The pair forms a two-sided certificate: the race
is reachable without the assumption and unreachable with it, so T2b is neither
vacuous nor decorative.

\medbfparagraph{T3 is not implied by T1 and T2}
It is natural to ask whether the claim obligation T3
(\code{IntermediaryMustClaim}) follows from the ordering constraints. It does
not. Removing T3 from \code{multihop.spthy} and re-running falsifies
\Cref{thm:intermediary}: \mtamarin\ returns a concrete 59-step counterexample.
T2b orders the two timeouts, $\mathit{TimedOut}(\mathit{ptrOut}) <
\mathit{TimedOut}(\mathit{ptrIn})$, but does not compel the node to
\emph{sweep} its incoming HTLC once the outgoing one is redeemed. That
obligation is precisely what T3 encodes, and it corresponds in deployment to
watchtower infrastructure~\cite{mccorry2019pisa}.

\begin{inv}[Assumption minimality]
\label{inv:minimality}
None of T1, T2a, T2b, T3 is implied by the others. T1--T2b carry timing
soundness; T3 alone carries intermediary safety.
\end{inv}

\medbfparagraph{Deriving what the core model assumes}
The concrete theory treats time relationally, as an ordering of events. The two
small natural-number theories discharge the arithmetic it takes for granted.
\code{Cltv.spthy} derives that the CLTV delta is strictly positive and that the
resulting claim window is non-empty across a staggered path.
\code{Clock.spthy} runs a live block clock through the full forward lifecycle
and derives the claim window between hops at concrete block heights
(\code{Claim\_Window\_Exists},
\code{Transitive\_Preimage\_Before\_Upstream\_Deadline}), together with the
lifecycle property that no HTLC may be added after a channel closes.

\subsection{Griefing and the Wormhole Attack}

Two pathologies survive an otherwise correct execution. Both are reachable in
our model, and we distinguish sharply between them.

\medbfparagraph{Griefing is an availability failure}
\code{Honest\_Intermediary\_Refunds} exhibits a trace in which capital is
locked and later refunded with no settlement anywhere on the path and no key
compromise. Every party follows the protocol; the receiver simply never
discloses $x$. No safety property is violated --- the funds return --- but every
forwarder's liquidity is immobilised until its deadline elapses. This is a
denial of service, not theft, and we report it as such.

\medbfparagraph{The wormhole, quantified}
The wormhole attack is known~\cite{malavolta2019anonymous} and we claim no
novelty for exhibiting it. Two non-adjacent colluding parties relay the
preimage out of band, bypassing the honest intermediary between them, which
forwards liquidity and earns nothing.
\code{Wormhole\_Fee\_Theft\_Reachable} finds this trace in 51 steps: the sender
settles and the receiver is paid, yet the middle hop is never redeemed and
neither intermediary records a \code{FeeEarned} event.

What is new is the quantification. Because our amounts are carried in the model
rather than abstracted away, we can state and prove not merely that a loss
occurs but exactly what it is.

\begin{thrm}[Wormhole theft is exactly the path fees]
\label{thm:wormhole}
There is a trace in which the sender locks $v + f_1 + f_2$, the receiver is
paid $v$, no key is compromised, and neither intermediary earns a fee.
\hfill(\code{Wormhole\_Steals\_Exactly\_The\_Fees}, 81 steps)
\end{thrm}

Prior symbolic treatments of agent-skipping establish that the honest
intermediary is bypassed, but their formalisms carry no numeric sort and so
cannot express \emph{how much} is lost. \Cref{thm:wormhole} pins the
adversary's gain to $f_1 + f_2$, the sum of the bypassed fees, which is the
economically meaningful statement.
